
\section*{Глава 3. Разработка веб-приложения}
\addcontentsline{toc}{section}{Глава 3. Разработка веб-приложения}
\stepcounter{section}

\subsection{Разработка базовой структуры клиентской части приложения}

Инициализация и сборка клиентской части веб-приложения была выполнена с помощью современного инструмента 
Vite в связке с библиотекой \textbf{React.js}. Данный выбор обеспечил высокую скорость перезагрузки 
модулей при разработке и оптимизацию итогового бандла. Для стилизации интерфейсов применен фреймворк Tailwind CSS, реализующий 
концепцию Utility-First, что позволило отказаться от громоздких каскадных таблиц стилей в пользу атомарных классов.

Архитектура клиентской части построена на основе компонентного подхода. Исходный код был логически разделен на следующие директории:
\begin{itemize}
    \item \texttt{components} — переиспользуемые элементы пользовательского интерфейса (шапка, подвал, карточка товара, боковые панели дашбордов);
    \item \texttt{pages} — уникальные страницы-компоненты (главная витрина товаров, страница авторизации, корзина, личный кабинет);
    \item \texttt{api} — модуль инкапсуляции сетевых запросов и настройки связи с серверной частью.
\end{itemize}

Маршрутизация в одностраничном приложении (SPA) реализована с помощью библиотеки \texttt{react-router-dom}. Был спроектирован единый диспетчер маршрутов в главном файле \texttt{App.jsx}, обеспечивающий переходы по страницам без перезагрузки браузера. 

Ключевым элементом базовой структуры является централизованный HTTP-клиент на базе библиотеки \textbf{Axios}. Для обеспечения безопасности и прозрачной работы с 
JWT-токенами был реализован механизм перехватчиков запросов. Данный архитектурный паттерн позволяет автоматически извлекать токен доступа из 
локального хранилища браузера (\texttt{localStorage}) и внедрять его в заголовки перед отправкой запроса на сервер. Использование централизованного 
HTTP-клиента исключило дублирование кода при обращении к эндпоинтам на сервере и заложило надежный фундамент для дальнейшей реализации модулей авторизации и 
управления контентом.

\subsection{Разработка серверной части, реализация аутентификации и авторизации}

В соответствии с задачей 3.2, серверная часть веб-приложения была разработана с использованием фреймворка \textbf{FastAPI}. Взаимодействие с реляционной системой 
управления базами данных PostgreSQL реализовано с помощью \textbf{SQLAlchemy ORM}, что позволило абстрагироваться от прямых SQL-запросов и обеспечить безопасность от инъекций.

Для организации надежной системы авторизации был выбран стандарт \textbf{OAuth2} с использованием JSON Web Tokens. Данный подход позволяет 
реализовать \textit{Stateless}-архитектуру, где сервер не хранит активные сессии пользователей в оперативной памяти. Это существенно повышает масштабируемость платформы. 

В БД была спроектирована модель пользователей (\texttt{users}) со связью к таблице ролей (\texttt{roles}). Пароли пользователей не сохраняются в 
открытом виде: перед записью в СУБД они проходят криптографическое хеширование с использованием алгоритма \texttt{bcrypt}. 

При успешной проверке учетных данных на эндпоинте \texttt{/auth/login} сервер генерирует подписанный JWT-токен, в полезную нагрузку которого внедряется
уникальный идентификатор пользователя.

Специфика разрабатываемого продукта требует строгого разделения прав доступа между покупателями, продавцами и администраторами. Для решения этой 
задачи в FastAPI был реализован параметризованный класс-зависимость \texttt{RoleChecker}. Такой механизм в проекте гарантирует, 
что валидация токена и проверка роли пользователя (например, `RoleChecker([3])`) будут выполнены до передачи управления бизнес-логике функции. В случае несовпадения 
прав доступа, система автоматически прервет обработку запроса и вернет клиенту статус-код \texttt{403 Forbidden}, обеспечивая надежную 
защиту платформы.

\subsection{Реализация CRUD-интерфейсов для пользователей, товаров и заказов}

В приложении был реализован комплекс модулей, обеспечивающих CRUD-операции для ключевых сущностей в системе. Логика управления данными была 
сегментирована в зависимости от ролевой принадлежности авторизованного пользователя.

Управление пользователями реализовано в рамках панели администратора. Был спроектирован интерфейс модерации, представляющий собой 
агрегированную таблицу с данными об атрибутах (логин, ФИО, контактные данные, текущая роль). Администратор наделен правами редактирования 
профилей для исправления ошибочных данных, изменения системных ролей, а также безвозвратного удаления учетных записей при нарушении правил 
платформы. Благодаря каскадному удалению на уровне базы данных, при удалении учетной записи автоматически очищаются связанные с ним товары и истории заказов.

Интерфейс управления каталогом товаров разделен между продавцами и администраторами. Для первой роли реализована изолированная B2B-панель. При создании и 
редактировании карточки товара продавец заполняет базовые поля и загружает графическое изображение. Загрузка медиафайлов реализована через 
отдельный независимый эндпоинт, возвращающий URL-адрес сохраненного на сервере изображения. 

Особое внимание при разработке формы товара было уделено работе с динамическими характеристиками, хранящимися в СУБД в типе данных \texttt{JSONB}. На клиентской 
части был спроектирован визуальный конструктор параметров. При вводе значений через запятую, алгоритм фронтенд-части автоматически преобразует 
строку в массив, что в дальнейшем позволяет системе корректно отображать эти параметры в виде кликабельных селекторов на витрине маркетплейса. 

Модуль управления заказами также адаптирован под специфику многопользовательской торговли:
\begin{itemize}
    \item \textbf{Для покупателя:} реализован интерфейс истории покупок в личном кабинете. Покупатель видит общую сумму чека, плановую дату доставки и 
	детализированный список купленных товаров с индивидуальными статусами их логистики.
    \item \textbf{Для продавца:} спроектирована панель обработки входящих заявок. Система автоматически фильтрует 
	глобальную таблицу позиций заказов, отображая продавцу только те товары, которые принадлежат его магазину. Продавец имеет возможность изменять 
	статус логистики через выпадающий список, что отражается в кабинете покупателя.
    \item \textbf{Для администратора:} реализован глобальный мониторинг всех транзакций платформы с информацией о покупателях  и 
	возможностью изменения статуса любых позиций (например, для разрешения спорных ситуаций).
\end{itemize}

Реализованный комплекс CRUD-операций полностью покрывает жизненный цикл сущностей приложения, обеспечивая надежное хранение данных и 
эргономичное взаимодействие пользователей с системой.

\subsection{Реализация системы корзины покупателя и бизнес-логики обработки заказов}

В рамках задачи был спроектирован и реализован модуль управления пользовательской корзиной и процесс оформления заказа. Для оптимизации нагрузки на 
серверную инфраструктуру и повышения отзывчивости интерфейса, промежуточное хранение данных корзины реализовано на стороне клиента с 
использованием \texttt{localStorage} браузера.

При нажатии на кнопку «Добавить в корзину» в карточке товара, клиентское приложение формирует объект, включающий идентификатор товара, количество и выбранные 
пользователем конфигурации (динамические селекторы цвета, памяти и т.д.). Для предотвращения дублирования позиций алгоритм проверяет корзину на наличие 
аналогичного товара с такими же выбранными параметрами. Сравнение объектов реализовано путем их сериализации в формат JSON. В случае совпадения увеличивается 
счетчик количества, в противном случае создается новая запись. Также на этапе добавления внедрена блокировка, не позволяющая превысить доступный складской остаток товара.

Интерфейс страницы корзины разделен на две логические области: список выбранных товаров с элементами управления количеством и панель итогового чека. При изменении 
количества товаров алгоритм динамически пересчитывает общую сумму и обновляет локальное хранилище. 

Бизнес-логика оформления заказа инициируется при заполнении пользователем адреса и плановой даты доставки. При отправке формы клиентское приложение передает на сервер 
сформированный массив позиций корзины. Обработка заказа на стороне бэкенда представляет собой сложную транзакцию, выполняющуюся в несколько этапов:
\begin{enumerate}
    \item \textbf{Проверка консистентности:} сервер повторно валидирует наличие запрошенных товаров в базе данных и проверяет, не превышает ли запрашиваемое количество текущие 
	складские остатки.
    \item \textbf{Списание остатков:} в случае успешной валидации система уменьшает значение поля \texttt{stock\_quantity} у соответствующих товаров, исключая 
	возможность их продажи другим пользователям.
    \item \textbf{Формирование чека:} алгоритм рассчитывает итоговую стоимость корзины на стороне сервера и создает запись заказа (\texttt{Order}).
    \item \textbf{Привязка позиций:} для фиксации связи «один-ко-многим», система использует метод \texttt{flush()} из ORM SQLAlchemy. Он принудительно отправляет 
	SQL-запрос в базу данных без закрытия транзакции, генерируя уникальный UUID заказа. Этот идентификатор присваивается всем дочерним позициям (\texttt{OrderItem}), после чего 
	транзакция окончательно фиксируется методом \texttt{commit()}.
\end{enumerate}

При успешном завершении транзакции локальная корзина клиента очищается, а заказ становится доступен в кабинете покупателя и панелях управления соответствующих продавцов.

\subsection{Разработка модуля отчетности и работы с файлами}

Был спроектирован и реализован модуль коммерческой аналитики, интегрированный в панель управления продавца. Целью данного модуля является предоставление агрегированных 
финансовых показателей и генерация отчетной документации.

Для просмотра статистических данныхна клиенте был разработан дашборд (рисунок ~\ref{fig:donut}), отображающий общую выручку, количество уникальных заказов и объем проданных 
единиц товара. Расчет метрик выполняется на стороне базы данных с применением агрегатных функций (подсчет уникальных чеков осуществляется через 
связку \texttt{func.count(func.distinct(...))}). Это переносит вычислительную нагрузку с приложения на СУБД, обеспечивая высокую скорость отклика.

\begin{figure}[H]
    \centering
    \includegraphics[width=1.0\textwidth]{donut.png}
    \caption{Раздел панели управления продавца с аналитической составляющей}
    \label{fig:donut}
\end{figure}

Визуализация структуры продаж реализована в виде динамической круговой диаграммы. На клиентской стороне алгоритм рассчитывает процентную долю каждой категории 
в общей выручке продавца и генерирует CSS-правило \texttt{conic-gradient}. Использованый подход позволил отказаться от использования сторонних библиотек для 
отрисовки графиков.

Особое внимание было уделено функционалу выгрузки финансовой отчетности в формате CSV. Пример экспортируемой таблицы на рисунке ~\ref{fig:exel_otchet}. Бизнес-логика экспорта 
реализована с учетом принципов оптимизации ввода-вывода и экономии дискового пространства сервера:
\begin{enumerate}
    \item \textbf{Сбор данных:} сервер выполняет SQL-запрос с объединениями (\textit{JOIN}) таблиц заказов, позиций и товаров, отфильтрованный по идентификатору текущего 
	продавца и отсортированный по убыванию даты транзакции.
    \item \textbf{Виртуализация файла:} вместо физического сохранения CSV-документа на жестком диске сервера, используется модуль \texttt{io.StringIO}, создающий текстовый буфер 
	непосредственно в оперативной памяти. 
    \item \textbf{Кодирование:} для корректного отображения кириллицы в начало потока принудительно записывается маркер последовательности 
	байтов \texttt{\textbackslash ufeff}. Запись строк данных осуществляется с применением разделителя «точка с запятой».
    \item \textbf{Подсчет итогов и стриминг:} в процессе итерации по списку продаж алгоритм рассчитывает итоговую выручку и дописывает ее в конец файла. Готовый буфер 
	кодируется в байты и отправляется клиенту через класс \texttt{StreamingResponse}. Потоковая передача данных защищает сервер от переполнения оперативной 
	памяти в случае генерации отчетов с сотнями тысяч строк.
\end{enumerate}

\begin{figure}[H]
    \centering
    \includegraphics[width=1.0\textwidth]{exel_otchet.png}
    \caption{Пример экспортируемой таблицы в формате csv}
    \label{fig:exel_otchet}
\end{figure}

Принятые архитектурные решения обеспечили высокую производительность модуля отчетности, гарантируя продавцам мгновенный доступ к актуальной аналитике и безопасный экспорт 
коммерческих данных без создания избыточной нагрузки на сервер.

Исходный код разработанного веб-приложения размещен в системе контроля версий GitHub\cite{github_code}.